\documentclass[11pt,a4paper]{scrartcl}

\usepackage[T1]{fontenc}
\usepackage[utf8]{inputenc}
\usepackage{kpfonts}
\usepackage[english]{babel}
\usepackage{amsmath,amssymb}
\usepackage[numbers,sort&compress]{natbib}
\usepackage[colorlinks=true,linkcolor=blue,citecolor=blue,urlcolor=blue]{hyperref}
\providecommand{\doi}[1]{\href{https://doi.org/#1}{doi: #1}}
\usepackage[margin=2.5cm]{geometry}
\usepackage{setspace}
\usepackage{enumitem}
\usepackage[automark,headsepline]{scrlayer-scrpage}

\onehalfspacing
\setkomafont{pagehead}{\small}

\title{Implication of HSAT2 Repeat Sequence RNAs in ME/CFS and Long COVID: A Hypothesis}

\author{Olivier Croce\textsuperscript{1}, Yannick Loth\textsuperscript{2},
        Patrick Lomonte\textsuperscript{3},\\[2pt]
        Andr{\'e} Trautman\textsuperscript{4},
        Genevi\`{e}ve Fourel\textsuperscript{1*}\\
       \small \textsuperscript{1} LBMC, ENS Lyon, CNRS, Inserm, Lyon, France\\
       \small \textsuperscript{2} Independent researcher, Luxembourg\\
       \small \textsuperscript{3} Universit\'e Claude Bernard Lyon 1, France\\
       \small \textsuperscript{4} Independent researcher, France\\[0.5em]
       \small *Corresponding author: \texttt{genevieve.fourel@ens-lyon.fr}}

\date{July 2026 \\ \vspace{0.5em} \small Prepared for the ISLC-PAIS Conference 2026}

\lohead{\small HSAT2 hypothesis in PAIS}
\rohead{\small \pagemark}

\begin{document}
\emergencystretch 1em
\maketitle

%───── Abstract (verbatim per Fourel) ─────
\begin{abstract}
  Human satellite~2 (HSAT2) is a pericentromeric repeat representing
  approximately 1.5\% of the human genome.  Under physiological
  conditions, it is transcriptionally silenced through DNA methylation
  and H3K9me3-dependent heterochromatin \cite{Shadle2019HSATII,
  Gerard2023CTCFpreprint}.  In cancer, however, HSAT2 may become
  aberrantly expressed and promotes immune suppression through
  extracellular vesicle-mediated reprogramming of myeloid cells,
  leading notably to the expansion of myeloid-derived suppressor cells
  and immune exhaustion \cite{Ruzanov2024ETSsatelliteEV}.

  Here we propose that an analogous mechanism may contribute to the
  development of post-acute infection syndromes (PAIS), including
  Long COVID and ME/CFS.  This model provides a potential explanation
  for the frequent observation that these disorders emerge following
  reactivation of latent herpesviruses.  Indeed, HSAT2 RNAs are also
  induced upon herpesvirus replication.  Herpesvirus infections are
  known to disrupt centromere organization and can induce transcription
  of pericentromeric satellite repeats, including HSAT2
  \cite{Lahaye2025centromere,Nogalski2019HSATII,Vourc_h2022HSF1SatIII}.
  Following repeated or closely spaced viral infections, cumulative
  changes in chromatin composition and loss of CpG methylation at
  pericentromeric loci may become sufficiently extensive to initiate a
  self-sustaining process of progressive epigenetic erosion
  \cite{BonnetFourel2026ProAB,deVega2014methylation,
  Helliwell2020methylation,Peppercorn2025methylation,
  Shadle2019HSATII}.

  By analogy with cancer, we hypothesize that HSAT2 RNAs released in
  extracellular vesicles could propagate this altered epigenetic state,
  driving persistent immune dysregulation and immune exhaustion
  \cite{Ruzanov2024ETSsatelliteEV}.  Rather than presenting definitive
  proof, this Perspective assembles the available evidence into a
  coherent mechanistic framework linking the initiating trigger to
  downstream molecular events and clinical manifestations.  It also
  generates specific, experimentally testable predictions that can now
  be tested using patient-derived samples
  \cite{Seimiya2023HSATII,Yoruker2026HSAT2cfDNA}.
\end{abstract}

%───── 1. HSAT2: a pericentromeric repeat, normally silenced ─────
\section{What HSAT2 Is and Why It Is Normally Silent}

HSAT2 (human satellite~2) is a stretch of DNA repeated many times
in the genome, occupying approximately 1.5\% of our DNA.  It
sits in pericentromeric regions: the DNA flanking each
chromosome's central pinch point (the centromere, where the two
arms meet).  Repetitive sequences account for roughly half of
the human genome---far more than the $\sim$2\% that encodes
proteins---yet have historically been neglected.

In healthy cells, HSAT2 does not produce RNA.  This silencing is
maintained by two locks: (i)~DNA methylation, in which small methyl
groups are added to specific sites (CpG, cytosine--phosphate--%
guanine dinucleotides), and (ii)~tightly packed chromatin called
heterochromatin, reinforced by a chemical mark on the histone proteins around which DNA is
wound (H3K9me3: histone~H3 trimethylated at lysine~9).  Together these prevent the cellular
machinery from transcribing HSAT2 into RNA
\cite{Shadle2019HSATII,Gerard2023CTCFpreprint}.

%───── 2. Cancer: HSAT2 → immune suppression ─────
\section{What Happens in Cancer}

In several cancers, HSAT2 escapes silencing and is transcribed
into RNA, which is packaged into extracellular vesicles
(EVs)---tiny membrane-wrapped particles, roughly 100~nm wide, that
cells shed into body fluids and travel via the bloodstream.
When HSAT2-laden EVs are taken up by immature cells of the myeloid
lineage (precursors of several innate immune cell types), they
reprogram those cells to dampen immunity.  Specifically, the EVs
drive expansion of myeloid-derived suppressor cells (MDSCs)---a
population of incompletely matured cells that suppress other
immune cells and expand in disease.

MDSCs suppress natural killer (NK) cells---rapid killers of
virus-infected and cancerous cells---and drive
CD8\textsuperscript{+} T-cells (``killer'' T-cells critical
for clearing infections) into exhaustion, a state of progressive
functional loss.  This is mediated by the
signalling molecules IL-10, IL-35, the enzyme IDO1, and TGF-$\beta$
(transforming growth factor beta)
\cite{Ruzanov2024ETSsatelliteEV}.  Tumours expressing HSAT2 are
immunologically cold: they show little immune-cell infiltration,
consistent with HSAT2 actively suppressing anti-tumour immunity
\cite{Ninomiya2023CAST}.

Critically, cells taking up HSAT2-containing EVs produce and release
their own HSAT2-carrying EVs, spreading the signal to additional
cells---a self-sustaining loop that can outlast the original cancer
trigger \cite{Ruzanov2024ETSsatelliteEV}.

%───── 3. HSAT2 in PAIS ─────
\section{Extension to Post-Acute Infection Syndromes}

We propose that this same chain of events may operate in post-acute
infection syndromes (PAIS)---most notably Long COVID and ME/CFS.
These conditions frequently begin after a viral illness, and many
patients experience reactivation of dormant herpesviruses (HSV-1,
the cold-sore virus; HCMV, cytomegalovirus; HHV-6).

The proposed sequence:

\textbf{Trigger.}  When a herpesvirus reactivates, it replicates at
the centromeres---the pinch points of chromosomes---where viral
proteins cause the local DNA to amplify abnormally
\cite{Lahaye2025centromere}.  This physical disruption destabilises
the adjacent pericentromeric regions and switches on HSAT2
transcription \cite{Nogalski2019HSATII,Vourc_h2022HSF1SatIII}.

\textbf{Cumulative erosion.}  Each reactivation damages chromatin
structure and strips the methylation marks keeping HSAT2 silent.
After repeated or closely spaced infections, cumulative loss can
cross a threshold: HSAT2 production becomes self-sustaining and
the epigenetic erosion continues without further viral stimulation
\cite{BonnetFourel2026ProAB,
deVega2014methylation,Helliwell2020methylation,
Peppercorn2025methylation,Shadle2019HSATII}.  DNA methylation
patterns are known to be altered in ME/CFS and Long COVID, and
pericentromeric repeats are especially susceptible to methylation
loss.

\textbf{Propagation.}  By analogy with cancer, HSAT2 RNAs from
cells experiencing centromere disruption are released into the
bloodstream in EVs, reach immature myeloid cells, drive MDSC
expansion, and induce T-cell exhaustion---contributing to the
NK-cell dysfunction and broader immune dysregulation documented in
ME/CFS and Long COVID \cite{Ruzanov2024ETSsatelliteEV}.  Because
EV-mediated propagation is self-sustaining, immune exhaustion could
persist long after the triggering infection has cleared.

%───── 4. Predictions ─────
\section{Testable Predictions}

The model generates four predictions testable in patient samples:

\begin{enumerate}[label=(\arabic*)]
  \item \textbf{EV HSAT2 elevation.}  HSAT2 RNA or DNA in
    blood-derived EVs will be elevated in a subgroup of ME/CFS
    and Long COVID patients with a confirmed viral trigger.
  \item \textbf{Patient subgroups.}  Patients who test positive for
    EV HSAT2 will show a distinct clinical profile: viral-onset
    history, expanded MDSC populations, and impaired NK-cell
    function.
  \item \textbf{Biomarker panels.}  A combined panel (EV HSAT2,
    MDSC frequency, cytokines) will identify the affected subgroup
    with high specificity.
  \item \textbf{Therapeutic response.}  In HSAT2-positive patients,
    methyl donors---compounds supplying the raw material for DNA
    methylation (such as SAMe [S-adenosylmethionine], methylfolate,
    and betaine)---will reduce EV HSAT2 more than placebo, whereas
    demethylating drugs (HDAC inhibitors, 5-azacitidine) would be
    predicted to worsen erosion.  \textit{These are testable
    predictions, not treatment recommendations.}
\end{enumerate}

These predictions can be pursued immediately.  Two HSAT2 detection
assays validated in cancer (TRAP-ddPCR for serum RNA and
hybridisation capture for plasma cell-free DNA) can be deployed
\cite{Seimiya2023HSATII,Yoruker2026HSAT2cfDNA}.  The simplest
starting point is RT-qPCR (reverse-transcription quantitative
PCR) on RNA from patient white blood cells (PBMCs) and
serum EVs, stratified by viral trigger.  Existing ME/CFS EV
sequencing data (Seifert et al.\ 2026 \cite{Seifert2026EVmiRNA})
could be re-analysed for HSAT2 without new samples.  Finding elevated HSAT2 in EVs from patients with
documented herpesvirus reactivation, co-varying with MDSC
expansion, would support the hypothesis; absence in a well-powered
cohort would argue strongly against it.

\textbf{Limitations.}  Each individual step is documented:
herpesviruses induce HSAT2 transcription
\cite{Nogalski2019HSATII}, HSAT2 RNA is packaged into EVs in cancer
\cite{Ruzanov2024ETSsatelliteEV}, and HSAT2 is detectable in human
serum \cite{Kishikawa2016SATII,Seimiya2023HSATII}.  However, the
chain connecting all three---viral centromere disruption leading to
HSAT2 in circulating EVs in a non-cancer context---has not been
demonstrated.  This Perspective is a framework for testing, not
established fact.

%───── Acknowledgments ─────
\section*{Acknowledgments}

We thank Claire Vourc'h (IAB, Universit\'e Grenoble Alpes) for
expertise on HSAT3---a related pericentromeric repeat---and for
initiating a task force to test the HSAT2/ME hypothesis, and
Martina Seifert (Charit\'e, Berlin) for discussion of
extracellular vesicle data.  This hypothesis was
developed by the authors.  Within an AI-assisted analysis of the
ME/CFS literature \cite{Loth2026MECFS}, Opus (Anthropic) evaluated
the hypothesis and identified strong mechanistic coherence and
cross-disease corroboration.  This builds on two decades of
research on repeat sequences as drivers of genome organisation.  Preprints
\cite{Gerard2023CTCFpreprint,BonnetFourel2026ProAB} are identified
as such.

%───── Bibliography ─────
\renewcommand{\doi}[1]{\href{https://doi.org/#1}{doi: #1}}
\bibliographystyle{unsrtnat}
\bibliography{refs}

\end{document}
